\section{Expectation as a continuous weighted average}

Expectation was introduced in the discrete case as a weighted average:
\[
\E[X]=\sum_x xP(X=x).
\]
Each possible value contributes its value multiplied by its probability mass. For
a continuous random variable, individual values usually have probability zero, so
this formula cannot be copied literally. The idea of a weighted average remains,
but the weights are supplied by density and accumulated by integration.

\begin{definitionbox}
Let \(X\) be a continuous random variable with density \(f_X\). If the integral is
well-defined, the \textbf{expectation} of \(X\) is
\[
\E[X]=\int_{-\infty}^{\infty}x f_X(x)\,dx.
\]
\end{definitionbox}

This formula should be read as a continuous weighted average. The value \(x\) is
weighted by the density near \(x\), and the integral adds all those infinitesimal
contributions.

\subsection*{Why the formula uses \(x f_X(x)\)}

In a small interval near \(x\) of width \(dx\), the probability is approximately
\[
f_X(x)\,dx.
\]
Values in that very small interval are approximately equal to \(x\). Therefore
the contribution of that small interval to the average is approximately
\[
x f_X(x)\,dx.
\]
Adding such contributions across all possible values gives
\[
\int_{-\infty}^{\infty}x f_X(x)\,dx.
\]
This is not a proof at the level of measure theory, but it gives the correct
interpretation: expectation is still an average, and density provides the
continuous weights.

\subsection*{Example: uniform distribution on \([0,1]\)}

Let \(X\) be uniformly distributed on \([0,1]\). Then
\[
f_X(x)=1
\]
for \(0\leq x\leq1\). Hence
\[
\E[X]=\int_0^1 x\,dx=\left[\frac{x^2}{2}\right]_0^1=\frac12.
\]
This agrees with the geometric intuition: the interval \([0,1]\) is symmetric,
and its center is \(1/2\).

\subsection*{Example: uniform distribution on \([a,b]\)}

If \(X\) is uniformly distributed on \([a,b]\), then
\[
f_X(x)=\frac1{b-a},\qquad a\leq x\leq b.
\]
Therefore
\[
\E[X]=\int_a^b x\frac1{b-a}\,dx
=\frac1{b-a}\left[\frac{x^2}{2}\right]_a^b.
\]
Thus
\[
\E[X]=\frac{b^2-a^2}{2(b-a)}=\frac{(b-a)(a+b)}{2(b-a)}=\frac{a+b}{2}.
\]
The expected value is the midpoint of the interval. This is not because the
midpoint is the most likely value; in a continuous uniform distribution no single
point has positive probability. It is the center of balance of the whole
interval.

\subsection*{Example: triangular density}

Let
\[
f_X(x)=
\begin{cases}
2x, & 0\leq x\leq1,\\
0, & \text{otherwise}.
\end{cases}
\]
Then
\[
\E[X]=\int_0^1 x\cdot 2x\,dx=\int_0^1 2x^2\,dx.
\]
Hence
\[
\E[X]=\left[\frac{2x^3}{3}\right]_0^1=\frac23.
\]
This expectation is larger than \(1/2\), which matches the shape of the density.
The density is higher near \(1\), so the distribution places more probability on
larger values.

\subsection*{Expectation of a function of a random variable}

Often we need the expectation of a transformed quantity, such as \(X^2\),
\(\sqrt X\), or \(e^X\). If \(g\) is a suitable function and \(X\) has density
\(f_X\), then
\[
\E[g(X)]=\int_{-\infty}^{\infty}g(x)f_X(x)\,dx,
\]
provided the integral is well-defined.

For example, if \(X\) is uniformly distributed on \([0,1]\), then
\[
\E[X^2]=\int_0^1 x^2\,dx=\frac13.
\]
This is different from \((\E[X])^2=(1/2)^2=1/4\). In general,
\[
\E[X^2]\neq (\E[X])^2.
\]
The first quantity averages squared values. The second squares the average.
These are different operations.

\subsection*{Linearity of expectation}

Linearity remains true in the continuous case. If the expectations exist, then
\[
\E[aX+bY]=a\E[X]+b\E[Y].
\]
This property does not require independence. It is a property of averaging, not
a property of separate random mechanisms.

For a single continuous random variable and constants \(a,b\), one can see this
from the integral formula:
\[
\E[aX+b]=\int_{-\infty}^{\infty}(ax+b)f_X(x)\,dx.
\]
Distributing the integral gives
\[
\E[aX+b]=a\int_{-\infty}^{\infty}x f_X(x)\,dx
+b\int_{-\infty}^{\infty}f_X(x)\,dx.
\]
Since the total integral of the density is \(1\),
\[
\E[aX+b]=a\E[X]+b.
\]

\subsection*{Expectation may not exist}

The integral defining expectation must be meaningful. Some distributions have
heavy tails, and the integral of \(|x|f_X(x)\) may be infinite. In such cases,
the expectation is not finite. This is not a technical detail to ignore: writing
down the formula is not enough; one must know whether the integral converges.

In this introductory course, most examples will have finite expectation. Still,
the condition is part of the mathematical definition.

\begin{summarybox}
For a continuous random variable, expectation is a continuous weighted average:
\[
\E[X]=\int_{-\infty}^{\infty}x f_X(x)\,dx.
\]
The density supplies the weights, and the integral accumulates them. The
interpretation is the same as in the discrete case, but sums are replaced by
integrals and point masses are replaced by density.
\end{summarybox}
